\songtitle{Mirá}

\tag*{1996}\tag{original-work}\tag{musical-theater}

\begin{notes}
Original work from 1996.\\
Musical theater piece using conversational narrative framing.
\end{notes}

\begin{style}
Musical theater showtune with folk-pop influences. %
Baritone male vocals deliver a narrative performance with clear diction and theatrical phrasing. %
The arrangement features a bright acoustic guitar strumming eighth-note patterns, a melodic accordion providing counter-melodies and sustained chords, and a double bass playing on beats 1 and 3. %
A light percussion section includes a tambourine on the backbeat and subtle shaker. %
The tempo is 115 BPM in 4/4 time, likely in the key of G major. %
The harmonic structure follows a diatonic progression typical of contemporary musical theater, with a clean, dry production style that emphasizes vocal clarity and acoustic instrument separation.
\end{style}

\begin{lyrics}

\songpart{Verse 1}
\annotation{strummed acoustic guitar, double bass, accordion}
Amiga, ¿puedo llamarte amiga? mirá! tengo un problema, tú lo conoces o lo has oído en parte, estoy seguro.  Resulta que, hace tiempo ya, encontré a una persona, tú la conoces, su nombre; si no lo sabes lo descubrirás. %
\annotation{tambourine enters on backbeat} Hace algún otro tiempo descubrí que la amaba, o al menos eso creí, o, ¿sabes?, creo que ni eso creí pero lo hice.  Y también hace algún tiempo, se lo hice saber.  No sé si me creyó 

\end{lyrics}

